\section{A transitive-kernel descent for a single standard layer}
\label{sec:standard-kernel-descent}

\paragraph{Status and shared provenance.}
HD-SS06, reserved question Q2ed, is a separate derivation by
\texttt{/root/g34\_eightfolds}. Root proposed the quotient-curve
inequality and the constant \(168\). The complete actual-cover,
field, primitive-line polarization and minimal-source interfaces
are written below. The already accepted HD-SS01--03 remains
unchanged. Complete author, root and independent whole-proof readings pass; the exact source and computation scopes are retained in their records.
No irreducible-induction conclusion is made.

Use the actual sphere setting of HD-SS01--02. Thus \(J\) is finite,
\(C\longrightarrow\mathbf P^1\) is a connected \(J\)-Galois cover,
and the faithful absolutely irreducible coefficient representation is
\[
 V=(\operatorname{Std}_n\circ x)\otimes U,\qquad
 x:J\longrightarrow S_n,\quad A_n\subset x(J),\qquad
 b=\dim U\ge2.
 \tag{SSK.1}
\]
Put \(a=n-1\). Both factors are actual finite linear representations. Assume that
\(V\) has a model over its exact imaginary quadratic character field
\(K\), and that its multiplicity variation has dimension six, Hodge
signature \((3,3)\), and actual period rank nine. The original primitive
factor \(B\) is selected by a rational Rosati-fixed primitive projector,
with the polarization restricted from the Jacobian. Put
\[
 K_0=\ker U,\quad H=x^{-1}((S_n)_1),\quad N=H\cap K_0,\quad
 X=C/H,\quad Y=C/N,\quad C_U=C/K_0,\quad Q=U(J).
 \tag{SSK.2}
\]
Here \(K_0\) is a subgroup, whereas \(K\) is the coefficient field.
The notation \(\mu_U\) and \(D'\) retains its precise meaning in
HD-SS02: \(D'\) is the moving-rank sum of
\(\tau=U|_H\) on the actual cover \(Y\longrightarrow X\).

\begin{theorem}[HD-SS06]
Suppose that \(K_0\) is transitive on the \(n\) letters and
\[
                              n>168(b+1).
 \tag{SSK.3}
\]
Then
\[
 \mu_U=0,\qquad F_U=F_\tau=K,\qquad
 U(H)=Q,\qquad g(X)=0,\qquad b\ge3.
 \tag{SSK.4}
\]
The actual connected \(Q\)-cover \(Y\longrightarrow X=\mathbf P^1\),
with faithful irreducible coefficient representation \(\tau\) of
degree \(b\), has an exact \(K\)-model and a primitive sixfold factor
\(A\). Its actual variation has period rank nine and signature
\((3,3)\). There is a rational \(K\)-isogeny \(\varphi:B\longrightarrow A\)
for which the natural polarizations satisfy
\[
 \varphi^*\lambda_A=c_0\lambda_B,\qquad c_0\in\mathbf Q_{>0},
 \qquad
 \delta(A)=\delta(B)
       \quad\hbox{in }\mathbf Q^\times/
                         \operatorname{Nm}_{K/\mathbf Q}(K^\times).
 \tag{SSK.5}
\]
Equality concerns the signed determinant norm class, not the integral
polarization type or the degree of the isogeny.

In particular, such a representation \(V\) cannot have the smallest
coefficient degree among actual exact-\(K\), rank-nine sources of
this natural polarization determinant type: the new degree \(b\)
is strictly smaller than \((n-1)b\). This is a source descent,
although \(\operatorname{Ind}_H^J\tau=U\oplus V\) remains reducible.
\end{theorem}

\begin{proof}
\paragraph{The actual quotient diagram and field equality.}
Transitivity of \(K_0\) gives \(J=K_0H\). Therefore
\[
 H/N\simeq J/K_0\simeq Q,\qquad [K_0:N]=[J:H]=n.
 \tag{SSK.6}
\]
In particular \(Y\longrightarrow X\) is a connected \(Q\)-Galois
cover and \(Y\longrightarrow C_U\) is an actual finite map of
degree \(n\). The latter need not be Galois. No map in the
opposite direction is assumed.

Here \(n\ge5\). The normal subgroup \(x(K_0)\) of \(x(J)\)
contains \(A_n\). Indeed, its intersection with the simple group
\(A_n\) is either trivial or all of \(A_n\). In the former case
it centralizes \(A_n\), since the commutator lies in that intersection.
The centralizer of \(A_n\) in \(S_n\) is trivial, contradicting
transitivity. Thus
\(K_0\cap x^{-1}(A_n)\) is an actual normal subgroup mapping onto
\(A_n\) and acting trivially on \(U\). It maps injectively as well:
an element in its kernel acts trivially on both factors of the
faithful \(V\).

The stronger actual-subgroup interface in HD-SS01 now applies.
For clarity, multiply any \(g\in J\) by an element of this subgroup
so that its permutation is either the identity or a transposition,
according to its parity. The \(U\)-matrix does not change, while
the standard trace becomes \(n-1\) or \(n-3\), both nonzero.
Consequently \(F_U=F_V=K\); HD-SS01 also gives \(F_\tau=F_U\)
and constructs \(K\)-models for \(U,\tau\) by multiplicity one.
The image equality \(U(H)=Q\) follows directly from \(J=K_0H\).
Thus \(\tau\) is faithful as a representation of \(Q\).

\paragraph{Positive partner cohomology forces a hyperbolic quotient.}
Suppose for a contradiction that \(\mu_U>0\). Since \(K_0\)
acts trivially on \(U\), quotient pullback identifies its multiplicity
in \(H^1(C,\mathbf C)\) with that in \(H^1(C_U,\mathbf C)\).
The latter cohomology vanishes in genus zero. In genus one its
holomorphic and antiholomorphic lines are each \(Q\)-invariant,
so every complex irreducible constituent has dimension one.
The irreducible \(U\) has dimension \(b\ge2\); hence
\(g(C_U)\ge2\).

The connected \(Q\)-Galois cover \(C_U\longrightarrow\mathbf P^1\)
therefore has positive orbifold excess
\[
 \Delta=\frac{2g(C_U)-2}{|Q|}
        =-2+\sum_i(1-1/d_i)\ge\frac1{42},
 \tag{SSK.7}
\]
where the \(d_i\ge2\) are its genuine inertia orders.
Only the elementary positive-excess calculation in the proof of
D-TD02, equation (TD.14), is used here, not that lemma's
nonsolvability hypothesis. Explicitly, at least five marks give
excess at least \(1/2\), and four give at least \(1/6\) unless
all orders are two, when the excess is zero. With three orders
\(u\le v\le w\), the smallest positive cases according to the
first two orders are \((2,3,7),(2,4,5),(3,3,4)\), or cases
with larger excess. Their minimum is \(1/42\).
At most two marks cannot give positive excess.

Riemann--Hurwitz for \(Y\longrightarrow C_U\), followed by the
\(Q\)-Galois cover \(Y\longrightarrow X\), gives, with \(h=g(X)\),
\[
 n\Delta
 \le\frac{2g(Y)-2}{|Q|}
 =2h-2+\sum_j(1-1/o_j)
 \le2h-2+D'
 =6+\mu_U-2(b-1)(h-1).
 \tag{SSK.8}
\]
Every genuine \(Q\)-inertia is nonidentity under the faithful
\(\tau\), so its moving rank is at least one and dominates
\(1-1/o_j\). This justifies the second inequality even when
different original marks induce different numbers of genuine marks.
The last equality is the exact HD-SS02 identity (SS.9).

The all-inertia estimate (SS.17), which keeps \(\mu_U\), says
\[
          \mu_U\le2b+\frac{4b+12}{n-2}<2b+1.
 \tag{SSK.9}
\]
The strict inequality follows from (SSK.3): it is stronger than
\(n>4b+14\). Hence the integer \(\mu_U\) is at most \(2b\).
Since \(h\ge0\) and \(b\ge2\), (SSK.8) is at most \(4b+4\).
But (SSK.3) and (SSK.7) give \(n\Delta>4b+4\).
This contradiction proves \(\mu_U=0\).
Period rank nine was not needed for this vanishing step.

\paragraph{The lower-degree actual sphere source.}
The exact variation identity (SS.7) now reduces to the \(V\)-variation.
All rational Galois conjugates of the inactive \(U\)-summand vanish,
since they occur in the rational cohomology of \(C\).
Thus \(\tau\) has multiplicity six, signature \((3,3)\) and actual
period rank nine. These statements concern the actual family
of \(Q\)-covers \(Y\longrightarrow X\), not just its Euler data.

For completeness the base must be rational. With \(\mu_U=0\),
\(D'=6-2b(h-1)\). If \(h=1\), there are at most six genuine
marks and the marked-base dimension is at most six. If \(h=2\),
the only possibilities with \(b\ge2\) are \(b=2,D'=2\) or
\(b=3,D'=0\), with moduli dimensions at most five or three.
For \(h\ge3\), \(D'<0\). Thus rank nine forces \(h=0\).
The connected sphere cover then has at least twelve genuine
marks and \(D'=2b+6\), whence \(b\ge3\).
The original family maps to a \(Q\)-Hurwitz component containing
this rank-nine variation; no assertion of submersivity of an
arbitrary forgetful map is required.

\paragraph{Actual polarized primitive factors.}
We supply the polarization argument rather than assume that an
arbitrary isogeny preserves its determinant. The exact \(K\)-models
and character fields identify the relevant rational group-algebra
components as \(\operatorname{Mat}_{ab}(K)\) for \(V\) and
\(\operatorname{Mat}_b(K)\) for \(\tau\).
Their Schur indices are one by the model construction in HD-SS01.
The multiplicity of \(\tau\) in \(V|_H\) is one.

Choose positive invariant \(K\)-Hermitian coefficient forms
\(h_V,h_\tau\), by finite averaging. The one-dimensional \(K\)-space
\(\operatorname{Hom}_H(\tau_K,V_K)\) has a nonzero map \(j\).
Schur's lemma and positivity give
\[
                    j^*h_V=c\,h_\tau,\qquad c\in\mathbf Q_{>0}.
 \tag{SSK.10}
\]
The forms are Hermitian over the imaginary quadratic field \(K\),
whose fixed field is \(\mathbf Q\); this is why the scalar is rational.
On the \(V\)-isotypic rational cohomology, the invariant polarization
has coefficient form \(h_V\) tensor its multiplicity Hermitian form.
This is the usual Schur decomposition of an invariant pairing.
Since \(\mu_U=0\), the entire active \(\tau\)-isotypic part of
\(H^1(C)\) comes from this multiplicity-one copy in \(V|_H\).
The other summand of \(\operatorname{Ind}\tau=U\oplus V\) contributes
zero cohomology.

To realize the comparison geometrically, choose a rational
Rosati-fixed primitive projector \(p_\tau\) in
\(\operatorname{Mat}_b(K)\subset\mathbf Q[Q]\). Such a projector
is the orthogonal projection onto a \(K\)-line for \(h_\tau\).
Use the canonical quotient-algebra embedding
\[
 \mathbf Q[Q]\longrightarrow e_N\mathbf Q[H]\subset\mathbf Q[J],
 \qquad e_N=\frac1{|N|}\sum_{u\in N}u .
 \tag{SSK.11}
\]
Its image \(p_H\) is self-adjoint. Rational Frobenius shows that
it has rank one over \(K\) in each of the \(U\)- and \(V\)-blocks
and is zero on every other irreducible rational block.
On \(H^1(C)\) the \(U\)-block is zero. Thus the connected image
\(p_HJ(C)\) is exactly a primitive \(V\)-factor.

Let \(A=p_\tau J(Y)\). For the actual quotient map
\(\pi:C\longrightarrow Y\), pullback restricts to an isogeny
from \(A\) onto \(p_HJ(C)\). The norm is its adjoint and
\(\pi_*\pi^*=[|N|]\), so the pullback of the restricted Jacobian
polarization is \(|N|\lambda_A\). Matrix units in
\(\operatorname{Mat}_{ab}(K)\) identify \(p_HJ(C)\) with the
chosen primitive \(B\) by a rational \(K\)-isogeny. All these maps
are rational linear combinations of deck correspondences, with
denominators cleared when actual homomorphisms are needed.
They work in families after a finite base change.
These are the actual correspondence and norm/pullback interfaces
of D-OG02, (OG.8)--(OG.9), and the final polarization argument
of D-BI08; the inactive extra summand has been removed explicitly here.

More explicitly, if \(t\) and \(v\) are nonzero coefficient vectors
on the chosen \(\tau\)- and \(V\)-lines, the ratio of the induced
multiplicity forms under this correspondence is
\[
             c_0=\frac{c\,h_\tau(t,t)}
                           {|N|\,h_V(v,v)}\in\mathbf Q_{>0}.
 \tag{SSK.12}
\]
One may invert the chosen rational isogeny; its direction only
inverts this positive rational ratio. Choosing the direction
\(B\longrightarrow A\) gives (SSK.5).
No hypothesis \(\operatorname{End}^0(B)=K\) is used: the isogeny
is the specified coefficient-line correspondence, not an arbitrary
isogeny with an uncontrolled multiplicity-space action.

With the same anti-invariant element of \(K\) used to normalize
both Hermitian forms, their rank is six. Hence
\[
                 c_0^6=\operatorname{Nm}_{K/\mathbf Q}(c_0^3).
 \tag{SSK.13}
\]
Changing a \(K\)-basis contributes another norm. Thus the signed
determinant classes agree. Positivity preserves the signature,
and the fixed correspondence preserves the actual rank-nine period
variation. Clearing a common denominator in this rational isogeny
multiplies \(c_0\) by a positive rational square, whose sixth power
is again a norm.

Finally \(b<(n-1)b\). Minimizing coefficient degree over the class
of actual exact-\(K\), rank-nine, naturally polarized sources of
the chosen determinant type excludes this configuration. The
comparison class includes the new connected \(Q\)-cover; it is
not restricted to representations irreducibly inducing \(V\).
This proves every assertion.
\end{proof}

\paragraph{Retained boundary.}
Without the size condition, or before the transitive-kernel
hypothesis is justified, no vanishing conclusion follows here.
In particular the real Euler-only example in HD-SS03 is not
being converted into an exact-\(K\) example, and no positive
\(\mu_U\) summand is discarded without the contradiction above.
